% BEGIN GENERATED ARTIFACT: prf:biconditional-equivalence-laws-propositional-logic
\newpage
\phantomsection
\label{prf:biconditional-equivalence-laws-propositional-logic}
\LRAProofFor{thm:biconditional-equivalence-laws-propositional-logic}

\begin{remark*}[Return]
\hyperref[thm:biconditional-equivalence-laws-propositional-logic]{Return to Theorem (Biconditional Equivalence Laws).}
\end{remark*}

\begin{theorem*}[Biconditional Equivalence Laws]
For all \(\varphi,\psi\in\WFF\),
\[
\varphi\leftrightarrow\psi
\equiv
(\varphi\to\psi)\land(\psi\to\varphi),
\]
and
\[
\varphi\leftrightarrow\psi
\equiv
(\varphi\land\psi)\lor(\neg\varphi\land\neg\psi).
\]
\end{theorem*}

\begin{proof}
\textbf{Professional Standard Proof.}~
TODO: supply the compact proof.
\end{proof}

\begin{proof}
\textbf{Detailed Learning Proof.}~
TODO: supply the detailed learning proof.
\end{proof}

\begin{remark*}[Proof structure]
Use the truth clause for biconditional, or rewrite using the conditional equivalence laws.
\end{remark*}

\begin{dependencies}
\begin{itemize}
  \item \hyperref[def:logical-equivalence-propositional-logic]{Logical Equivalence}
  \item \hyperref[thm:conditional-equivalence-laws-propositional-logic]{Conditional Equivalence Laws}
\end{itemize}
\end{dependencies}

\clearpage
% END GENERATED ARTIFACT: prf:biconditional-equivalence-laws-propositional-logic
